\section{Direct On-Policy Distillation}
\label{sec:method}

% \subsection{Motivation}
% \label{sec:motivation}

% We now formalize the policy-shift signal introduced above.
% Standard on-policy distillation (OPD) copies the post-RL teacher's final
% distribution, which mixes the useful change induced by RL with the weak
% teacher's capacity limits. \method{} keeps the same on-policy interface but
% changes the transferred object: it uses the teacher/reference policy shift as an
% improvement signal for the student, rather than using the weak teacher's final
% policy as an imitation target.

\subsection{Preliminaries: On-Policy Distillation}
\label{sec:preliminaries}

We consider autoregressive language-model policies. Let $x\sim\mathcal{D}$ be a
prompt and $y=(y_1,\ldots,y_T)$ a response, so that a policy $\pi$ factorizes as
$\log\pi(y\mid x)=\sum_{t=1}^{T}\log\pi(y_t\mid s_t)$ over prefixes
$s_t=(x,y_{<t})$. We distinguish four policies: the policy being trained
$\pi_\theta$ and its initial checkpoint $\pi_S$; and a teacher pair consisting of
a pre-RL reference $\pi_{T_{\mathrm{ref}}}$ and the corresponding post-RL teacher
$\pi_T$. Our weak-to-strong setting is the regime in which $\pi_T$ may be smaller or weaker than the target student, yet the transition from $\pi_{T_{\mathrm{ref}}}$ to $\pi_T$ still encodes information worth transferring.

On-policy distillation (OPD) supervises the student on trajectories it samples
itself. Given a prompt $x$, the student draws $\hat y\sim\pi_\theta(\cdot\mid x)$,
and at each visited prefix $s_t=(x,\hat y_{<t})$ we read off the student and
teacher next-token distributions $p_t(v)=\pi_\theta(v\mid s_t)$ and
$q_t(v)=\pi_T(v\mid s_t)$. Standard OPD minimizes the sequence-level KL to the
teacher on these rollouts,
\begin{equation}
\mathcal{L}_{\mathrm{OPD}}(\theta)
=\mathbb{E}_{x\sim\mathcal{D}}\big[\,
D_{\mathrm{KL}}\big(\pi_\theta(\cdot\mid x)\,\big\|\,\pi_T(\cdot\mid x)\big)\,\big],
\end{equation}
which factorizes exactly into dense, per-token supervision on the states the
student actually visits,
\begin{equation}
\mathcal{L}_{\mathrm{OPD}}(\theta)
=\mathbb{E}_{x\sim\mathcal{D},\,\hat y\sim\pi_\theta}\Big[\,
\textstyle\sum_{t=1}^{T}D_{\mathrm{KL}}(p_t\,\|\,q_t)\,\Big].
\end{equation}
In practice the teacher is queried only on the student's top-$k$ support
$S_t=\operatorname{TopK}_v p_t$, giving the local estimator
\begin{equation}
\mathcal{L}^{\text{top-}k}_{\mathrm{OPD}}(\theta)
=\mathbb{E}_{x\sim\mathcal{D},\,\hat y\sim\pi_\theta}\Big[\,
\textstyle\sum_{t=1}^{T}
D_{\mathrm{KL}}\big(\bar p_t^{S_t}\,\big\|\,\bar q_t^{S_t}\big)\,\Big],
\end{equation}
where $\bar p_t^{S_t}$ and $\bar q_t^{S_t}$ renormalize $p_t$ and $q_t$ on
$S_t$~\citep{li2026rethinkingopd}. This top-$k$ KL estimator produces a dense
per-token training signal involving the teacher--student log-ratio
$\log q_t(v)-\log p_t(v)$. \method{} inherits this on-policy, top-$k$ interface unchanged, and departs from OPD only in \textit{what signal is read from the teacher} 
(Section~\ref{sec:reward}).

\subsection{Policy Shifts as Implicit Rewards}
\label{sec:reward}

\textbf{Which signal to transfer.} Rather than imitating the teacher's absolute output distribution, \method{} transfers its RL-induced \textbf{policy shift}: the sequence-level log-ratio between the post-RL teacher $\pi_T$ and its pre-RL reference $\pi_{T_{\mathrm{ref}}}$,
\begin{equation}
    \Delta_T(y\mid x) = \log \pi_T(y\mid x) - \log \pi_{T_{\mathrm{ref}}}(y\mid x),
    \label{eq:dopd-delta}
\end{equation}
which isolates the teacher's RL-induced \textit{direction} rather than its
absolute final distribution. This shift is not an arbitrary heuristic: under
the \textbf{policy-as-reward} view of KL-regularized RL, it can be interpreted
as the teacher's implicit reward.

For a reward $r$, reference policy $\pi_{\mathrm{ref}}$, and penalty $\beta>0$,
the KL-regularized objective
$\max_\pi\mathbb{E}_{y\sim\pi}\!\big[r(x,y)-\beta\log\tfrac{\pi(y\mid x)}{\pi_{\mathrm{ref}}(y\mid x)}\big]$
admits the closed-form optimum $\pi^*\propto\pi_{\mathrm{ref}}\exp(r/\beta)$, hence
\begin{equation}
\log\frac{\pi^*(y\mid x)}{\pi_{\mathrm{ref}}(y\mid x)}
=\tfrac{1}{\beta}\,r(x,y)-\log Z(x).
\end{equation}
Since $\log Z(x)$ is constant across responses to the same prompt, the
policy--reference log-ratio recovers the reward up to a positive scale and a
per-prompt constant. This is the same identity that underlies Direct Preference
Optimization, which fits a policy from preferences without a separately trained
reward model~\citep{rafailov2023direct}. Here we use the identity in the
opposite direction: we are already \textit{given} a post-RL teacher $\pi_T$ and
its pre-RL reference $\pi_{T_{\mathrm{ref}}}$, and we read the reward-like signal
back out of their policy ratio. Treating $\pi_T$ as the KL-regularized optimum of some latent reward $r_T$ anchored at $\pi_{T_{\mathrm{ref}}}$,
\begin{equation}
    \Delta_T(y\mid x) = \tfrac{1}{\beta}\, r_T(x,y) - \log Z_T(x),
    \label{eq:dopd-implicit-reward}
\end{equation}
so the shift $\Delta_T$ can be interpreted as the teacher's \textbf{implicit
reward}, up to a positive scale and a per-prompt constant. \method{} transfers
this reward-like signal to the student.

\subsection{The Direct-OPD Objective}
\label{sec:dopd}

\textbf{Idealized sequence-level objective.}
We first write \method{} as a sequence-level objective. With the student
initialized from $\pi_S$, \method{} optimizes the student against the
teacher-shift signal $\Delta_T$ while regularizing the student toward its own
initialization,
\begin{equation}
    J_{\text{\method}}(\theta) = \mathbb{E}_{x\sim\mathcal{D}} \left[
    \mathbb{E}_{y\sim\pi_\theta(\cdot\mid x)} \big[ \Delta_T(y\mid x) \big]
    - \alpha\, D_{\mathrm{KL}} \big( \pi_\theta(\cdot\mid x) \,\|\, \pi_S(\cdot\mid x) \big) \right],
    \label{eq:dopd-obj}
\end{equation}
where $\alpha>0$ controls the KL penalty. In this idealized form, the objective
is itself a KL-regularized RL objective on the student: its optimum is
\begin{equation}
    \pi^*(y\mid x) \propto \pi_S(y\mid x)\, \exp\!\big(\tfrac{1}{\alpha}\Delta_T(y\mid x)\big)
    = \pi_S(y\mid x) \left( \frac{\pi_T(y\mid x)}{\pi_{T_{\mathrm{ref}}}(y\mid x)} \right)^{1/\alpha},
    \label{eq:dopd-opt}
\end{equation}
and substituting the implicit-reward form of Eq.~\ref{eq:dopd-implicit-reward}
gives $\pi^*\propto\pi_S\exp(r_T/\alpha\beta)$. The student is therefore
optimized as if it were running KL-regularized RL with the small teacher's
implicit reward, while using its own initialization $\pi_S$ as the reference.
Crucially, it obtains this reward-like signal entirely from the checkpoint pair, without querying a verifiable reward or running sparse-reward RL on the target.

\textbf{From sequence to token level.}
Because both teachers factorize autoregressively over the same prefixes, the
sequence shift decomposes exactly into a sum of token-level shifts,
\begin{equation}
    \Delta_T(y\mid x) = \sum_{t} r_t(y_t\mid s_t),
    \qquad
    r_t(v) = \log \pi_T(v\mid s_t) - \log \pi_{T_{\mathrm{ref}}}(v\mid s_t),
    \label{eq:dopd-token-reward}
\end{equation}
so that $r_t(v)$ can be viewed as a \emph{dense immediate per-token reward}:
$r_t(v)>0$ where the teacher's RL encouraged token $v$, and $r_t(v)<0$ where it
suppressed it. In practice, we use the corresponding zero-discount token-level
surrogate for Eq.~\ref{eq:dopd-obj}, crediting each candidate token by its
immediate shift $r_t(v)$ rather than by a future return over later shifts. This
reuses the on-policy, top-$k$ interface of Section~\ref{sec:preliminaries}.

\textbf{Top-$k$ action-space restriction.}
The implementation uses a local top-$k$ approximation to this zero-discount
token-level surrogate. To keep the signal on actions the student actually
considers, at each visited prefix we restrict to its top-$k$ support
$S_t=\operatorname{TopK}_v \pi_\theta(v\mid s_t)$ and renormalize the student
probabilities on it,
\begin{equation}
    \bar p_t(v) =
    \frac{\pi_\theta(v\mid s_t)}
    {\sum_{u\in S_t}\pi_\theta(u\mid s_t)},
    \qquad v\in S_t,
\end{equation}
which is the renormalized student distribution $\bar p_t^{S_t}$ of
Section~\ref{sec:preliminaries}.

\textbf{Analytical top-$k$ policy gradient.}
The immediate per-token reward $r_t(v)$ depends only on the teacher pair, so we
use it directly as the local advantage. A single-sample, token-level policy
gradient for the zero-discount surrogate would weight the log-likelihood of the
\emph{sampled} token by its immediate reward,
\begin{equation}
    \nabla_\theta J_{\mathrm{MC}} =
    \mathbb{E}_{x, y} \Big[
    \textstyle\sum_t r_t(y_t)\,
    \nabla_\theta \log \pi_\theta(y_t \mid s_t)
    \Big],
\end{equation}
but estimating each step from one token is high-variance. Since we already have
the full top-$k$ rewards and probabilities at every visited prefix, we replace
the single-token estimate at each step by its expectation over the restricted
distribution $\bar p_t$, i.e. we Rao--Blackwellize the per-step action while
still sampling the trajectory $y\sim\pi_\theta$:
\begin{equation}
    \nabla_\theta J_{\mathrm{analytical}} =
    \mathbb{E}_{x,\, y\sim\pi_\theta} \Big[
    \sum_t \sum_{v \in S_t}
    \underbrace{\textcolor{blue}{\bar p_t(v)}}_{\text{weight}} \,
    \underbrace{\textcolor{red}{r_t(v)}}_{\text{reward}} \,
    \nabla_\theta
    \underbrace{\textcolor{magenta}{\log \pi_\theta(v \mid s_t)}}_{\text{log-likelihood}}
    \Big].
\end{equation}
This gives a Rao--Blackwellized per-step estimator under the restricted top-$k$
action distribution; it removes the token-sampling variance of
$\nabla_\theta J_{\mathrm{MC}}$ while leaving the sampled trajectory distribution
untouched.

\textbf{Stop-gradient coefficient.}
For this surrogate to match the policy-gradient form, the weight $\bar p_t(v)$
must enter only as a scalar. It depends on $\theta$ through the student softmax,
so differentiating it would inject extra Jacobian terms from the top-$k$
distribution. We therefore detach the weighted reward into a static coefficient,
\begin{equation}
    A_t^{\mathrm{w}}(v) =
    \operatorname{stop\_gradient}\!\big(
    \textcolor{blue}{\bar p_t(v)} \cdot
    \textcolor{red}{r_t(v)}
    \big),
\end{equation}
giving the final local top-$k$ surrogate for the zero-discount objective,
\begin{equation}
    \nabla_\theta J_{\text{\method}} \approx
    \mathbb{E}_{x \sim \mathcal{D},\, y \sim \pi_\theta} \Big[
    \sum_t \sum_{v \in S_t}
    A_t^{\mathrm{w}}(v)\,
    \nabla_\theta
    \textcolor{magenta}{\log \pi_\theta(v \mid s_t)}
    \Big]
    - \alpha\, \nabla_\theta
    D_{\mathrm{KL}}\big(
    \pi_\theta(\cdot\mid x) \,\|\, \pi_S(\cdot\mid x)
    \big).
\end{equation}
In practice we do not differentiate the KL analytically; we use the standard
KL-penalty implementation in \texttt{verl} to anchor $\pi_\theta$ to $\pi_S$.

\subsection{Adaptive KL Control}
\label{sec:method-adaptive-kl}

The reward \method{} transfers carries a scale it does not get to choose. By
Eq.~\ref{eq:dopd-implicit-reward}, $\Delta_T=\tfrac{1}{\beta}r_T-\log Z_T$: its magnitude is the teacher's reward divided by the teacher's KL budget $\beta$, both fixed when the teacher was trained and encoded in how far its RL pushed $\pi_T$ from $\pi_{T_{\mathrm{ref}}}$, yet neither is recoverable from the checkpoint pair.
The per-token shift $r_t(v)=\log\pi_T(v\mid s_t)-\log\pi_{T_{\mathrm{ref}}}
(v\mid s_t)$ is likewise a function of the teacher pair alone; the student enters
only through \emph{where} the signal is read off---the prefixes its rollouts
visit and the top-$k$ actions retained at each one.

This scale mismatch makes a \emph{fixed} $\alpha$ unlikely to transfer robustly
across settings. From the optimum in Eq.~\ref{eq:dopd-opt},
$\pi^*\propto\pi_S\exp(\tfrac{1}{\alpha}\Delta_T)$, $\alpha$ is the temperature
of an exponential tilt of $\pi_S$, and is meaningful only relative to the scale
of $\Delta_T$. But $\Delta_T$ lives in the teacher's reward units
($\tfrac1\beta r_T$) while $\alpha$ is a coefficient on the \emph{student's}
KL---two scales with no a priori conversion, the teacher's being unobservable.
A single $\alpha$ therefore cannot be calibrated in advance across teacher--student pairs, and we instead adopt a lightweight controller that adapts $\alpha$ to the running scale of the dense signal.

Let $\alpha_m$ be the KL coefficient at training iteration $m$, and let
$\bar r_m$ be the batch-level mean of the student-weighted shift
$\bar p_t(v)\,r_t(v)$ over the visited prefixes and their top-$k$ candidates---the
same dense reward the gradient optimizes. Before the actor update we set
\begin{equation}
    \alpha_{m+1}
    =
    \operatorname{clip}
    \left(
    \alpha_m
    \left(1 + \epsilon\,\operatorname{sgn}(\bar r_m)\right),
    \alpha_{\min},
    \alpha_{\max}
    \right),
    \label{eq:adaptive-kl}
\end{equation}
with $\operatorname{sgn}(0)=0$; by default $\epsilon=0.01$ and
$[\alpha_{\min},\alpha_{\max}]=[0.5,2.5]$. Here $\operatorname{sgn}(\bar r_m)$
reports whether, on the prefixes the student currently visits, the teacher's RL
on average raised or lowered the retained candidate tokens---a batch-level sign
that the student-weighted shift carries directly.

The controller is intended to keep the dense signal balanced. When the visited
prefixes carry positive teacher shift on average, it \emph{raises} $\alpha$ to
discourage over-amplifying that local signal; when the average shift is
negative, it \emph{lowers} $\alpha$, weakening the anchor so the dense gradient
can move probability away from tokens whose likelihood the teacher's RL reduced.
This differs from the standard adaptive controller for an in-reward KL penalty,
which steers toward a target KL value: here the update is driven by the
\emph{sign of the \method{} dense reward} and acts on the explicit actor KL loss
coefficient. We analyze KL-coefficient sensitivity in
Section~\ref{sec:kl-coefficient}.

% A central practical finding is that different teacher--student pairs require
% different KL coefficients to reach their best performance. This sensitivity is
% expected in \method{}: the scale of the teacher/reference log-ratio depends on the
% teacher pair, the student, and the region of prefixes visited by the current
% policy. We therefore use a simple adaptive controller for the actor KL loss
% coefficient $\alpha$, which sets how far the student may move from its initial
% policy.

% Let $\alpha_m$ be the KL coefficient at training iteration $m$, and let
% $\bar{A}_m$ be the batch-level mean of the valid candidate advantages
% $A_t^{\text{w}}(v)$. Before the actor update, we set
% \begin{equation}
%     \alpha_{m+1}
%     =
%     \operatorname{clip}
%     \left(
%     \alpha_m
%     \left(1 + \epsilon\,\operatorname{sgn}(\bar{A}_m)\right),
%     \alpha_{\min},
%     \alpha_{\max}
%     \right),
% \end{equation}
% with $\operatorname{sgn}(0)=0$. In the implementation, the default $\epsilon=0.01$,
% and the coefficient is clipped to $[\alpha_{\min},\alpha_{\max}]=[0.5,2.5]$.

% This controller targets the balance point of the dense transfer signal rather
% than a fixed KL divergence. If the sampled prefixes have positive teacher-shift
% reward on average, the KL coefficient is increased to prevent the student from
% over-amplifying that local signal; if the average reward is negative, the
% coefficient is decreased, weakening the reference anchor so the dense reward
% gradient can move the student away from tokens whose probability became lower
% after teacher RL. This differs from the standard adaptive KL controller used for
% an in-reward KL penalty: here the update is driven by the sign of the \method{}
% dense reward and controls the explicit actor KL loss coefficient. We provide a
% detailed analysis of KL-coefficient sensitivity in
% Section~\ref{sec:kl-coefficient}.
